\documentclass[11pt]{letter}

\usepackage[margin=0.5in]{geometry}
\usepackage{hyperref}

\begin{document}

\begin{letter}{}

\opening{Dear Hiring Team,}

I am writing to apply for the Data Scientist position, Job Requisition ID 86144, with Tax and Revenue Administration at the Government of Alberta. With a Master of Science in Computer Science from the University of Calgary and extensive applied experience in machine learning, deep learning, large-scale data processing, statistical analysis, and research communication, I am excited about the opportunity to support evidence-based decision-making within Treasury Board and Finance.

My graduate research focused on developing scalable machine learning and data-management solutions for large geospatial and Earth observation datasets. I designed, trained, and evaluated deep neural networks in Python and PyTorch, performed exploratory data analysis and feature engineering, and built reproducible pipelines for data cleaning, transformation, model training, validation, and performance analysis. My work involved datasets containing more than one million candidate data patches and required careful attention to data quality, sampling strategy, generalization, scalability, and computational efficiency.

One of my primary research projects examined the reliability and generalizability of deep-learning models for digital elevation model super-resolution. I developed data-processing workflows using Python, NumPy, pandas, and PyTorch; implemented convolutional neural-network architectures and customized loss functions; and evaluated model performance using statistical and terrain-specific metrics. Through systematic experimentation, I identified how random data splitting could produce overly optimistic results and designed region-based validation methods that provided a more accurate assessment of real-world performance. This experience demonstrates my ability to apply systems thinking, recognize broader dependencies within analytical workflows, and develop practical solutions to complex data problems.

My second major research area involved developing a hierarchical framework for managing and retrieving large volumes of spatiotemporal satellite data. This work required integrating data from multiple sources, designing multiresolution data-processing methods, assessing storage and retrieval performance, and communicating technical results through peer-reviewed publications and presentations. These projects strengthened my understanding of data architecture, ETL-style workflows, reproducibility, and the challenges involved in making large datasets usable for operational analysis.

I also bring experience communicating complex technical concepts to diverse audiences. As a teaching assistant at the University of Calgary, I supported students in data science, programming, database, and analytical coursework. I explained technical methods, evaluated analytical work, provided structured feedback, and helped students translate data into defensible conclusions. These responsibilities developed my ability to communicate clearly with both technical and non-technical stakeholders and align analytical work with practical objectives.

My experience aligns closely with the competencies required for this position. I have demonstrated creative problem solving by designing new modeling and evaluation approaches when conventional methods were insufficient. I have shown agility by adapting research plans, data pipelines, and technical methods as requirements and evidence evolved. I have maintained a strong drive for results by independently managing complex projects from data acquisition and experimentation through publication and presentation. I also understand the importance of responsible data use, clear documentation, model transparency, privacy, and reproducible analytical practices, particularly when working with sensitive or decision-relevant information.

In addition to my core machine learning background, I have experience with SQL, statistical analysis, data visualization, neural networks, Git-based development, high-performance computing environments, and cloud-related concepts. I am also actively expanding my knowledge of MLOps, model deployment, cloud platforms, large language models, embeddings, retrieval-augmented generation, and responsible AI. My research background has trained me to learn emerging technologies quickly and apply them carefully to real organizational problems.

I am particularly interested in this opportunity because it combines advanced analytics with public service. The opportunity to help Tax and Revenue Administration improve reporting, program effectiveness, operational decision-making, and responsible use of data is highly meaningful to me. I would welcome the opportunity to contribute my analytical skills, research experience, and commitment to reliable and transparent data science to the Government of Alberta. I am available to work full-time in person in Edmonton and am prepared to participate in the required assessment and security-screening processes.

Thank you for considering my application. I look forward to the opportunity to discuss how my experience can contribute to the Business Intelligence and Analytics team within Tax and Revenue Administration.

\closing{Sincerely,

Amir Mirzai Golpayegani

\href{mailto:amirgolpaa24@gmail.com}{amirgolpaa24@gmail.com}}

\end{letter}

\end{document}